\documentclass[fontsize=15pt]{jlreq}

%プリアンブル
\usepackage{amsmath}
\pagestyle{empty}
\usepackage[top=15mm, bottom=20mm, left=20mm, right=20mm]{geometry}

%本文
\begin{document}
\begin{center}
{\LARGE {振り返り}} \\
-生物\textcircled{\scriptsize 1}- %単元ごとに変えるか消す
\end{center}
\vspace{1em}
\begin{flushright}
名前：\underline{\hspace{5cm}}
\end{flushright}

% 問題部分
\begin{enumerate}
    \item 細胞の核の中に見られるひものようなものをなんというか．\\
    ( \hspace{2cm} )
    \vspace{2em} 

    \item 生物の形質を決めるものを何と言うか．\\
    ( \hspace{2cm} )
    \vspace{2em} 

    \item 細胞分裂と体細胞分裂についてそれぞれ説明せよ．
    \begin{itemize}
      \item 細胞分裂：
      \vspace{2em}
      \item 体細胞分裂：
    \end{itemize}
    \vspace{2em} 

    \item たまねぎの根は，
     先端に近い部分の細胞の大きさが，
    根元に近い部分と比べて
    比較的小さいことがわかる．
    この理由を，体細胞分裂の観点から説明せよ． \\
    (\hspace{15cm})
    \vspace{2em} 

    \item タマネギの根の細胞を観察するために，
    はじめうすい塩酸に浸し，
    1分間温め，水の中でゆすいだ．
    なぜこのような操作が必要であるか．
    (\hspace{15cm})
    \vspace{2em} 

\end{enumerate}

\end{document}
